\documentclass[11pt,a4paper]{article}
\usepackage{fontspec}
\setmainfont{texgyreheros-regular.otf}[BoldFont=texgyreheros-bold.otf]
\usepackage[polish]{babel}
\usepackage{amsmath,amssymb,geometry,tabularx,array,fancyhdr}
\geometry{left=18mm,right=18mm,top=17mm,bottom=18mm}
\setlength{\parindent}{0pt}
\setlength{\parskip}{5pt}
\setlength{\tabcolsep}{5pt}
\renewcommand{\arraystretch}{1.15}
\newcolumntype{Y}{>{\raggedright\arraybackslash}X}
\pagestyle{fancy}\fancyhf{}
\fancyfoot[L]{\footnotesize Ćwiczenie 71 \quad Karta pomiarów i obliczeń}
\fancyfoot[R]{\footnotesize\thepage\ / 25}
\renewcommand{\headrulewidth}{0pt}
\newcommand{\titlepagepart}[1]{{\Large\bfseries Dyfrakcja światła}\par\smallskip{\large\bfseries #1}\par\medskip}
\newcommand{\note}[1]{{\small #1}\par}
\newcommand{\writebox}[2]{\begin{tabularx}{\linewidth}{|Y|}\hline\textbf{#1}\\\hline\rule{0pt}{#2}\\\hline\end{tabularx}\par\smallskip}
\begin{document}

\titlepagepart{1 Dane stanowiska}
\note{Zapisz własne odczyty. Niepewności możesz obliczyć później, ale dane o przyrządach zbierz na zajęciach.}
\begin{tabularx}{\linewidth}{|p{72mm}|Y|}\hline
\textbf{Co wpisać} & \textbf{Miejsce na wpis}\\\hline
\rule{0pt}{9mm}Imię i nazwisko / współwykonawcy & \\\hline
\rule{0pt}{9mm}Grupa i zespół / data pomiarów & \\\hline
\end{tabularx}\par\smallskip
\begin{tabularx}{\linewidth}{|p{82mm}|Y|Y|}\hline
\textbf{Wielkość i znaczenie} & \textbf{Wartość} & \textbf{Niepewność standardowa}\\\hline
\rule{0pt}{10mm}$\lambda\ [\mathrm{nm}]$ -- długość fali &  & \\\hline
\rule{0pt}{10mm}$L_2\ [\mathrm{mm}]$ -- szczeliny podwójne do detektora &  & \\\hline
\rule{0pt}{10mm}$R_2\ [\Omega]$ -- opór w serii &  & \\\hline
\rule{0pt}{10mm}$q_p\ [\mathrm{mm}]$ -- rozdzielczość położenia &  & \\\hline
\rule{0pt}{10mm}$q_U\ [\mathrm{mV}]$ -- rozdzielczość miernika &  & \\\hline
\rule{0pt}{10mm}$U_{\mathrm{bg}}\ [\mathrm{mV}]$ -- tło, opcjonalnie &  & \\\hline
\end{tabularx}\par\smallskip
\begin{tabularx}{\linewidth}{|p{72mm}|Y|}\hline
\textbf{Co wpisać} & \textbf{Miejsce na wpis}\\\hline
\rule{0pt}{10mm}Model miernika i specyfikacja dokładności & \\\hline
\rule{0pt}{10mm}Oznaczenie płytek / wartości nominalne & \\\hline
\rule{0pt}{10mm}Źródło informacji o długości fali & \\\hline
\end{tabularx}\par\smallskip
\writebox{Szkic układu: laser, szczelina, detektor; zaznacz odległość $L$}{25mm}
\note{Pomiar dla dwóch szczelin: krok $0{,}1\,\mathrm{mm}$. Długość fali podana w instrukcji: $\lambda=650\,\mathrm{nm}$. Potwierdź parametry aktualnego stanowiska.}
\clearpage
\titlepagepart{2 Ściąga z oznaczeniami}
\note{Tę stronę czytasz przed pomiarami. Wpisy i obliczenia wykonujesz w tabelach na dalszych stronach.}
\begin{tabularx}{\linewidth}{|p{43mm}|Y|}\hline
\textbf{Symbol} & \textbf{Co oznacza}\\\hline
\rule{0pt}{10mm}$p\ [\mathrm{mm}]$ & Surowy odczyt położenia detektora ze skali.\\\hline
\rule{0pt}{10mm}$I\ [\mathrm{j.u.}]$ & Odczyt miernika proporcjonalny do jasności, np. napięcie w mV. Jedną skalę zachowaj w całej serii.\\\hline
\rule{0pt}{10mm}$p_0,\ I_0$ & Położenie maksimum centralnego i sygnał w tym punkcie.\\\hline
\rule{0pt}{10mm}$x=p-p_0$ & Odległość od środka obrazu; nie jest tym samym co surowy odczyt skali.\\\hline
\rule{0pt}{10mm}$p_l,\ p_p$ & Położenia lewego i prawego ekstremum tego samego rzędu.\\\hline
\rule{0pt}{10mm}$r=\dfrac{p_p-p_l}{2}$ & Połowa odległości między parą ekstremów.\\\hline
\rule{0pt}{10mm}$m=1,2,\ldots$ & Numer minimum lub prążka po jednej stronie. Dla pary maksimów $-2,+2$ wpisujesz $m=2$.\\\hline
\rule{0pt}{10mm}$d\ [\mathrm{mm}]$ & Odległość między środkami dwóch szczelin.\\\hline
\rule{0pt}{10mm}$u(z)$ & Niepewność standardowa wielkości $z$, w tej samej jednostce co $z$.\\\hline
\end{tabularx}\par\smallskip
\note{Na stanowisku: ustaw obraz, dobierz opór bez przekroczenia zakresu miernika i nie zmieniaj go w serii. Dwie szczeliny: obejmij maksimum centralne i kilkanaście prążków interferencyjnych. Zapisz $L$, $\lambda$ i dokładności przyrządów.}
\writebox{Ustalenia prowadzącego i własne uwagi}{35mm}
\clearpage
\titlepagepart{Podwójna szczelina -- karta 1 z 20}
\note{Wpisuj $p$ ze skali oraz $I$ z miernika. W uwagach zaznacz maksimum, minimum lub zmianę warunków.}
\begin{tabularx}{\linewidth}{|Y|Y|Y|Y|}\hline
\textbf{Seria / płytka} & \textbf{$L_2\ [\mathrm{mm}]$} & \textbf{$R_2\ [\Omega]$} & \textbf{Krok $\Delta p\ [\mathrm{mm}]$}\\\hline
\rule{0pt}{7mm} &  &  & $0{,}1$ \\\hline
\end{tabularx}\par\smallskip
\begin{tabularx}{\linewidth}{|Y|Y|Y|Y|}\hline
\textbf{$p_0\ [\mathrm{mm}]$} & \textbf{$I_0\ [\mathrm{j.u.}]$} & \textbf{Skala sygnału, np. mV} & \textbf{Tło, opcjonalnie}\\\hline
\rule{0pt}{7mm} &  &  & $0{,}1$ \\\hline
\end{tabularx}\par\smallskip
\begin{tabularx}{\linewidth}{|p{9mm}|Y|Y|Y|}\hline
\textbf{Nr} & \textbf{Położenie $p\ [\mathrm{mm}]$} & \textbf{Sygnał $I\ [\mathrm{j.u.}]$} & \textbf{Uwagi}\\\hline
\rule{0pt}{7mm}1 &  &  & \\\hline
\rule{0pt}{7mm}2 &  &  & \\\hline
\rule{0pt}{7mm}3 &  &  & \\\hline
\rule{0pt}{7mm}4 &  &  & \\\hline
\rule{0pt}{7mm}5 &  &  & \\\hline
\rule{0pt}{7mm}6 &  &  & \\\hline
\rule{0pt}{7mm}7 &  &  & \\\hline
\rule{0pt}{7mm}8 &  &  & \\\hline
\rule{0pt}{7mm}9 &  &  & \\\hline
\rule{0pt}{7mm}10 &  &  & \\\hline
\rule{0pt}{7mm}11 &  &  & \\\hline
\rule{0pt}{7mm}12 &  &  & \\\hline
\rule{0pt}{7mm}13 &  &  & \\\hline
\rule{0pt}{7mm}14 &  &  & \\\hline
\rule{0pt}{7mm}15 &  &  & \\\hline
\rule{0pt}{7mm}16 &  &  & \\\hline
\rule{0pt}{7mm}17 &  &  & \\\hline
\rule{0pt}{7mm}18 &  &  & \\\hline
\rule{0pt}{7mm}19 &  &  & \\\hline
\rule{0pt}{7mm}20 &  &  & \\\hline
\end{tabularx}\par\smallskip
\note{Pomiary 1--20 z 400. Krok przesuwania detektora: $0{,}1\,\mathrm{mm}$. Zakres pomiaru musi obejmować wszystkie wymagane prążki.}
\clearpage
\titlepagepart{Podwójna szczelina -- karta 2 z 20}\titlepagepart{23 Podwójna szczelina -- obliczenia}
\note{Na wykresie liniowym ponumeruj jasne maksima, przypisując zero maksimum głównemu. Dobierz pary $-m,+m$.}
\[r_m=\frac{p_{p,m}-p_{l,m}}2,\qquad d_m=\frac{m\lambda L_2}{r_m}.\]
\begin{tabularx}{\linewidth}{|Y|Y|Y|Y|Y|}\hline
\textbf{$m$} & \textbf{$p_{l,m}\ [\mathrm{mm}]$} & \textbf{$p_{p,m}\ [\mathrm{mm}]$} & \textbf{$r_m\ [\mathrm{mm}]$} & \textbf{$d_m\ [\mathrm{mm}]$}\\\hline
\rule{0pt}{10mm}1 &  &  &  & \\\hline
\rule{0pt}{10mm}2 &  &  &  & \\\hline
\rule{0pt}{10mm}3 &  &  &  & \\\hline
\rule{0pt}{10mm}4 &  &  &  & \\\hline
\rule{0pt}{10mm}5 &  &  &  & \\\hline
\rule{0pt}{10mm}6 &  &  &  & \\\hline
\rule{0pt}{10mm}7 &  &  &  & \\\hline
\end{tabularx}\par\smallskip
\begin{tabularx}{\linewidth}{|Y|Y|Y|}\hline
\textbf{$N$} & \textbf{$\bar d=\dfrac1N\sum_m d_m\ [\mathrm{mm}]$} & \textbf{$u_{\mathrm A}(\bar d)\ [\mathrm{mm}]$}\\\hline
\rule{0pt}{12mm} &  & \\\hline
\end{tabularx}\par\smallskip
\[Q=\frac{I_{\min}}{I_{\max}}.\]
\begin{tabularx}{\linewidth}{|Y|Y|Y|}\hline
\textbf{Odczyt do oceny jakości} & \textbf{$p\ [\mathrm{mm}]$} & \textbf{$I\ [\mathrm{j.u.}]$}\\\hline
\rule{0pt}{10mm}Największe maksimum &  & \\\hline
\rule{0pt}{10mm}Najbliższe minimum &  & \\\hline
\end{tabularx}\par\smallskip
\begin{tabularx}{\linewidth}{|Y|Y|}\hline
\textbf{$Q\ [1]$} & \textbf{Sygnał surowy czy po korekcie tła?}\\\hline
\rule{0pt}{10mm} & \\\hline
\end{tabularx}\par\smallskip
\note{$Q$ bliskie 0: wyraźne prążki; $Q$ bliskie 1: prążki zanikają.}
\writebox{Uwagi o symetrii, obwiedni i jakości obrazu}{23mm}
\clearpage
\titlepagepart{24 Niepewności i końcowe wyniki}
\note{Wpisz niepewności standardowe i uzasadnij ich przyjęcie. Błąd graniczny oraz rozdzielczość nie są automatycznie niepewnością standardową.}
\[u_{\mathrm B}=\frac{\Delta}{\sqrt3}\quad(\text{błąd jednostajny }\pm\Delta),\qquad u_{\mathrm B}=\frac{q}{\sqrt{12}}\quad(\text{zaokrąglenie o kroku }q).\]
\begin{tabularx}{\linewidth}{|Y|Y|Y|}\hline
\textbf{Wielkość} & \textbf{Przyrząd / źródło i model} & \textbf{$u$ z jednostką}\\\hline
\rule{0pt}{11mm}$\lambda$ &  & \\\hline
\rule{0pt}{11mm}$L_2$ &  & \\\hline
\rule{0pt}{11mm}$p_l,p_p$ &  & \\\hline
\rule{0pt}{11mm}$I$ lub $U$ &  & \\\hline
\end{tabularx}\par\smallskip
\[
u_{\mathrm A}(\bar z)=\sqrt{\frac{\sum_{i=1}^N(z_i-\bar z)^2}{N(N-1)}},\qquad
u(r)=\frac12\sqrt{u^2(p_l)+u^2(p_p)}.
\]
\[
\frac{u(z)}{z}=\sqrt{\left(\frac{u(\lambda)}{\lambda}\right)^2+
\left(\frac{u(L)}{L}\right)^2+\left(\frac{u(r)}r\right)^2},\quad z=d_m.
\]
\note{Wzór względny dotyczy pojedynczego oszacowania przy niezależnych danych wejściowych. W średniej uwzględnij wspólne $\lambda,L$, korelacje. Sama $u_{\mathrm A}$ nie jest pełną niepewnością.}
\writebox{Obliczenia i uzasadnienie końcowej niepewności}{27mm}
\begin{tabularx}{\linewidth}{|Y|Y|Y|Y|Y|}\hline
\textbf{Wynik} & \textbf{Wartość [mm]} & \textbf{$u$ [mm]} & \textbf{$100u/z$ [\%]} & \textbf{$U=ku$ [mm]}\\\hline
\rule{0pt}{12mm}$\bar d$ &  &  &  & \\\hline
\end{tabularx}\par\smallskip
\begin{tabularx}{\linewidth}{|Y|Y|}\hline
\textbf{Współczynnik rozszerzenia $k$, jeśli stosujesz $U$} & \textbf{Wartość}\\\hline
\rule{0pt}{9mm} & \\\hline
\end{tabularx}\par\smallskip
\clearpage
\titlepagepart{25 Notatki do sprawozdania}
\note{Wpisuj obserwacje wynikające z własnych danych. Na osi logarytmicznej pokazuj tylko dodatnie sygnały i opisz pominięte punkty.}
\writebox{Cel, rzeczywista metoda i ustawienie aparatury}{24mm}
\writebox{Wykres dla dwóch szczelin: natężenie w funkcji położenia detektora}{24mm}
\writebox{Zgodność z teorią: minima, maksima, symetria, natężenia i wskaźnik $Q$}{24mm}
\writebox{Źródła rozbieżności oraz ich wpływ na wynik}{24mm}
\writebox{Wnioski na podstawie pomiarów}{24mm}
\note{Podstawa: instrukcja 71\_wykon.pdf oraz opis 71\_opis (1).pdf. Przed oddaniem sprawdź jednostki, niepewności, podpisy wykresów i zaokrąglenia.}
\end{document}
